% TOP_PROBLEM: {"id": 30007106, "title": "Sharp terminal edge count in random triangle removal", "category_id": 2, "category": {"id": 2, "name": "combinatorics", "display_name": "Combinatorics", "description": "Counting problems, graph theory, discrete structures.", "slug": "combinatorics", "order_index": 2, "created_at": "2026-07-31T15:26:25.671Z"}, "status": "open"}
% FIRST_POSTED: 2026-10-06
% ATTEMPT_STATUS: solved
% OPENAI_CLAIM: {"schema_version":1,"source_repo":"https://github.com/openai/math","source_commit":"adc7f1241b42e322a6451854ab7e4b4c146bf78a","release_date":"2026-10-06","match":"full","resolution":"proved","independently_reviewed":false,"second_pass":"agreed","families":[{"family":"188","title":"The sharp terminal leave in random triangle removal","subject":"Combinatorics","manuscripts":[{"title":"The sharp terminal leave in random triangle removal","pdf_path":"preprints/The-Sharp-Terminal-Leave-in-Random-Triangle-Removal-September-25-2026/The-Sharp-Terminal-Leave-in-Random-Triangle-Removal-September-25-2026.pdf","readme_path":"preprints/The-Sharp-Terminal-Leave-in-Random-Triangle-Removal-September-25-2026/README.md","date":"2026-09-25","theorem_ref":"preprints/The-Sharp-Terminal-Leave-in-Random-Triangle-Removal-September-25-2026/build/introduction.tex#thm:main"}],"lean":{"doc_path":"lean/docs/188.md","comparators":["lean/ComparatorChallenges/TriangleRemoval.json","lean/ComparatorChallenges/TriangleRemoval.lean"],"declarations":[],"covers_main_theorem":true},"reasoning_trace":null}]}
\documentclass[11pt]{article}
\usepackage[margin=25mm]{geometry}
\usepackage{fontspec}
\setmainfont{Latin Modern Roman}
\usepackage{amsmath,amssymb,mathtools}
\usepackage{unicode-math}
\setmathfont{Latin Modern Math}
\usepackage{enumitem}
\usepackage{xurl,hyperref,bookmark}
\hypersetup{colorlinks=true,urlcolor=blue}
\setcounter{secnumdepth}{0}
\setlength{\parindent}{0pt}
\setlength{\parskip}{5pt}
\setlength{\emergencystretch}{3em}
\begin{document}
\section*{Sharp terminal edge count in random triangle removal}

\subsection{OpenAI's claim}
\textit{Random triangle-removal terminal constant}: The final edge count in uniform random triangle removal from $K_n$, divided by $n^{3/2}$, is claimed to approach $1/(2\sqrt2)$ in mean square. Convergence of expectation and in probability follow.

\subsection{Outline of the argument}
\textit{Random triangle-removal terminal constant}: The process is stopped while subgraph counts remain controlled, then its continuation is represented by independent triangle priorities. Backward recursive survival queries are unfolded over the growing terminal time interval. Stability and paired-query estimates have errors smaller than the two-root survival scale, permitting first- and second-moment computations with the sharp constant.

\subsection{Scope relative to this problem}
\textit{Random triangle-removal terminal constant}: The manuscript’s result is exactly the terminal-edge-count mean-square limit E[(F\_n/n\textasciicircum{}(3/2)-1/(2sqrt(2)))\textasciicircum{}2]->0 for uniform triangle selection until no triangles remain. It concerns the terminal leave rather than an intermediate-time spread estimate.

With the hypotheses and conventions above, the claimed result addresses the whole mathematical assertion on this page.

\subsection{Formal verification}
Family 188 describes Lean coverage for the terminal triangle-removal mean-square limit with its sharp coefficient.

\textit{Random triangle-removal terminal constant}: The manuscript links \texttt{TriangleRemoval.lean}.

The documented selection includes this manuscript’s main assertion; that assertion retains the restrictions stated above.

This site has not run the Lean development. These links describe OpenAI’s selected formal statements and should be read with their stated scope.

\subsection{Status caveat}
OpenAI reports that an internal model produced the manuscripts. Their arguments have not been independently verified as proofs by this site. Any portion outside the documented formal scope may contain mistakes, and the catalogue label reports the announced mathematical claim.

\subsection{Sources}
\begin{itemize}
\item OpenAI, \emph{The sharp terminal leave in random triangle removal}, 2026-09-25; citation key \texttt{OAI:The-Sharp-Terminal-Leave-in-Random-Triangle-Removal-September-25-2026}. \href{https://github.com/openai/math/blob/adc7f1241b42e322a6451854ab7e4b4c146bf78a/preprints/The-Sharp-Terminal-Leave-in-Random-Triangle-Removal-September-25-2026/The-Sharp-Terminal-Leave-in-Random-Triangle-Removal-September-25-2026.pdf}{PDF (GitHub)}; \href{https://github.com/openai/math/raw/adc7f1241b42e322a6451854ab7e4b4c146bf78a/preprints/The-Sharp-Terminal-Leave-in-Random-Triangle-Removal-September-25-2026/The-Sharp-Terminal-Leave-in-Random-Triangle-Removal-September-25-2026.pdf}{PDF (direct)}; \href{https://github.com/openai/math/blob/adc7f1241b42e322a6451854ab7e4b4c146bf78a/preprints/The-Sharp-Terminal-Leave-in-Random-Triangle-Removal-September-25-2026/README.md}{Manuscript page and citation}; \href{https://github.com/openai/math/blob/main/preprints/The-Sharp-Terminal-Leave-in-Random-Triangle-Removal-September-25-2026/The-Sharp-Terminal-Leave-in-Random-Triangle-Removal-September-25-2026.pdf}{Latest version}.
\item Family 188: \href{https://github.com/openai/math/blob/adc7f1241b42e322a6451854ab7e4b4c146bf78a/lean/docs/188.md}{Lean scope}.
\item \href{https://github.com/openai/math/blob/adc7f1241b42e322a6451854ab7e4b4c146bf78a/lean/ComparatorChallenges/TriangleRemoval.json}{Comparator TriangleRemoval.json}.
\item \href{https://github.com/openai/math/blob/adc7f1241b42e322a6451854ab7e4b4c146bf78a/lean/ComparatorChallenges/TriangleRemoval.lean}{Comparator TriangleRemoval.lean}.
\item \href{https://github.com/openai/math/blob/adc7f1241b42e322a6451854ab7e4b4c146bf78a/CONTENTS.md}{OpenAI catalogue entry}.
\end{itemize}
\end{document}
